\documentclass{article}
\usepackage{graphicx}
\usepackage{subcaption}
\begin{document}
\begin{figure}
    \centering
    \begin{subfigure}[t]{0.4\linewidth}
        \includegraphics[width=\linewidth]{manuscript/figures/carnivores-kappa}
        \caption{Estimated per-branch kappa against the age of each branch. Each point is the posterior mean of the kappa parameter for one branch, shown on a logarithmic scale; the vertical bars show its 95\% credible interval and the horizontal bars show the 95\% credible interval of the age. Age is the height of the branch's midpoint above the present, measured in expected substitutions per site (the tree is estimated in units of substitutions and is not calibrated to time).}
        \label{fig:carnivoreskappa}
    \end{subfigure}
    \begin{subfigure}[t]{0.4\linewidth}
        \includegraphics[width=\linewidth]{manuscript/figures/carnivores-model-trees}
        \caption{Internal node age estimates under the two models. Each point is one internal node, with its age under the base (single-kappa) model on the horizontal axis and its age under the per-branch kappa variation model on the vertical axis; points are posterior means, bars show the 95\% posterior quantile interval, and the dotted line is the identity. Points falling below the identity for older nodes show that the kappa variation model infers younger ages for deeper nodes.}
        \label{fig:carnivorestree}
    \end{subfigure}
    
\end{figure}
\end{document}
